% =====================================================================
% PARTE I — Fundamentação                       [~5 min]
% Fonte: Chapters/fundamentacao
% =====================================================================

\section{Fundamentação}

\begin{frame}{Séries temporais: representação e pressupostos}
  \begin{itemize}
    \item Três casos de observação: amostragem uniforme, sequência ordenada
          por índice e série de eventos.
    \item O que a tese \alerta{não} assume.
      \pendente{listar as delimitações do Cap. 2}
    \item Estacionariedade, regime e deriva: distinção operacional em
          relação a anomalia.
  \end{itemize}
\end{frame}

\begin{frame}{Sinal, ruído, resíduo e anomalia}
  \begin{block}{Distinção central}
    Resíduo elevado \alerta{não implica} anomalia.
  \end{block}
  \vspace{0.4em}
  \begin{itemize}
    \item Decomposição aditiva e multiplicativa.
      \pendente{trazer eq:aditivo e eq:multiplicativo do Cap. 2}
    \item Estatísticas estimadas somente em partições autorizadas.
  \end{itemize}
\end{frame}

\begin{frame}{Qualidade de dados: ausência, corrupção e integridade}
  \begin{itemize}
    \item Mecanismos de ausência: \dest{MCAR}, \dest{MAR}, \dest{MNAR}.
    \item Implicação sobre a validade da imputação.
    \item \pendente{ligar aos requisitos derivados no Apêndice B}
  \end{itemize}
\end{frame}
